\chapter{Introdução e objetivos}

\section{Contexto}

O projeto \textbf{MiniLang / MiniLexer} foi desenvolvido para a disciplina de Linguagens Formais e Autômatos com o objetivo de relacionar os conceitos de Expressões Regulares e autômatos finitos com uma aplicação computacional funcional. A aplicação implementa um analisador léxico de uma linguagem de programação simplificada denominada MiniLang.

O enunciado da atividade exige uma aplicação com entrada, processamento e saída, no mínimo cinco Expressões Regulares relevantes, tratamento de entradas inválidas, código organizado e, para cada expressão, a apresentação do alfabeto, linguagem reconhecida, ER formal, sintaxe implementada, explicação dos operadores, AFN$\varepsilon$ e testes com cadeias aceitas e rejeitadas. Também estabelece que a ER formal, a sintaxe usada no programa, os testes e o AFN$\varepsilon$ devem representar a mesma linguagem.

\section{Objetivo geral}

Construir e validar um analisador léxico da MiniLang, demonstrando uma cadeia completa de consistência:

\begin{center}
\texttt{especificação $\rightarrow$ ER formal $\rightarrow$ AFN$\varepsilon$ $\rightarrow$ regex Python $\rightarrow$ testes $\rightarrow$ lexer}
\end{center}

\section{Objetivos específicos}

\begin{itemize}
    \item definir seis Expressões Regulares próprias e relevantes para a MiniLang;
    \item construir os AFN$\varepsilon$ correspondentes pela convenção de Thompson adotada no projeto;
    \item implementar um simulador genérico de AFN$\varepsilon$;
    \item implementar o MiniLexer com prioridade e maior lexema;
    \item validar a correspondência entre AFN$\varepsilon$ e regex Python por testes automatizados;
    \item validar estruturalmente estados, estados finais, alfabetos e transições dos seis AFN$\varepsilon$;
    \item disponibilizar os diagramas em JFLAP e Mermaid;
    \item documentar limitações, resultados e possibilidades de evolução.
\end{itemize}

\section{Escopo}

A MiniLang é tratada exclusivamente no nível léxico. O projeto não implementa parser, análise semântica, compilação, execução, verificação completa de tipos, geração de código ou máquina virtual. Esta delimitação evita que funcionalidades posteriores alterem o foco acadêmico do trabalho.
